\section{Constructing the local indices and their canonical replacement}
\label{sec:canonical-indices}

\status{CI1--CI18 reconstruct the finite local calculations left open in
TL20: finite admitted cosets, lattice normal forms, preservation of first
unfolding nonvanishing, descent of the actual projected character, the
auxiliary-prime one-coset calculation, and the resulting initial and
canonical local-index records. CI19 identifies the exact contribution
to the construction of \texttt{TC.Ok}. Arithmetic host, character, tuple,
and ambient theta data remain hypotheses where stated; no analytic
theta identity or stage of \texttt{Later} is inferred from these local
calculations. These are mathematical reconstructions from the pinned
source, without a Lean build or kernel-verification claim.}

\paragraph{Local conventions.}
Let $P$ be a nonzero prime of the integers of the number field $E$,
$O=\mathcal O_{E,P}$, and $q=|\mathcal O_E/P|$. Except in the finite
coset argument CI3, assume the literal \texttt{DegOneUnram} condition:
$q$ is a rational prime and $q\notin P^2$. Thus $q$ is a uniformizer
of the discrete valuation ring $O$, and $O/qO=\mathbb F_q$. The
coefficient field $K$ has characteristic zero. All local lattices below
are in rational $E$-vector spaces, rather than completions.

At step $j$, put $n_j=k+j(t+1)$ and $c_j=m-1-j$, for
$0\le j\le m-1$. The source rank assumptions include $m\ge4$ even
and $k\ge m+3$. Write
\[
 Y_j=X_j\oplus\bigoplus_{a=0}^{c_j-1}Ee_a,
 \qquad X_j=E^{n_j}.
 \tag{CI.1}
\]
The group $R_j$ preserves $X_j+Ee_0+\cdots+Ee_{a-1}$,
fixes $e_0$, and sends $e_a$ to $e_a$ plus an earlier vector.
Its action on $X_j$ may be any invertible linear map; calling
all of $R_j$ unipotent would be incorrect. Its kernel on $X_j$ is
$N_j$. Set
\[
 a(r)=\sum_{a=0}^{c_j-2}(r e_{a+1})_{e_a}.
 \tag{CI.2}
\]
Then $a(rs)=a(r)+a(s)$, as follows by multiplying these triangular
matrices. At a twist prime the basic lattice is $O^{n_j+c_j}$;
at an auxiliary, or type, prime it is
$\Lambda_j=O^{n_j}\oplus\bigoplus_a q^{a+1}Oe_a$.
The words ``nonzero unfolding'' always mean the source predicate
\texttt{FiniteSumNonzeroLargeDepth}: all sufficiently deep box sums
are nonzero as functions of the retained matrix variable. It is not
an assertion about an unproved analytic integral.

\paragraph{CI1: finite quotients, stable box sums, and affine transport.}
For every $s\ge1$, the integers $0,\ldots,q^s-1$ represent $O/q^sO$
exactly once. Indeed they represent $O/qO$; after removing that first
digit and dividing by $q$, induction gives existence. Uniqueness
follows because the $P$-valuation of a nonzero rational integer is
its usual $q$-adic valuation. Consequently
\[
 \mathcal B_{A,B}=\{0,1,\ldots,q^{A+B}-1\}/q^A
 \tag{CI.3}
\]
represents $q^{-A}O/q^BO$.

Suppose $f:E^d\to K$ is supported in $q^{-A_0}O^d$ and is invariant
under $q^{B_0}O^d$. For $A\ge A_0$, $B\ge B_0$, define
\[
 I(f)=q^{-Bd}\sum_{u\in\mathcal B_{A,B}^d}f(u).
 \tag{CI.4}
\]
Increasing $A$ only adds zero support cosets; increasing $B$ by $b$
splits each old period coset into $q^{bd}$ equal terms. The factor
$q^{-Bd}$ cancels this multiplicity. Thus (CI.4) is independent of
both sufficiently large depths. The same proof applies uniformly
to a family with common support and period bounds.

For an invertible $E$-linear map $T$ and $b\in E^d$, the functions
$f$ and $f(Tu+b)$ have bounded support and a full period lattice.
Their stable sums vanish simultaneously. To see this without a
measure theorem, choose an $O$-lattice $L$ which is a period lattice
for both functions. The change of variable identifies the sum over
$L$-cosets for $f(Tu+b)$ with the sum over $TL$-cosets for $f$.
Use the common sublattice $H=L\cap TL$, which has finite index in
each, to compare either sum with the sum over $H$-cosets. The two
comparisons multiply by the positive integers $[L:H]$ and $[TL:H]$.
Every box normalization is also a nonzero rational number.
Characteristic zero therefore proves the asserted equivalence.
Bounds on the entries of $T,T^{-1},b$ choose the needed box depths;
there is no unstated convergence argument.

\paragraph{CI2: a nonzero finite sum forces an affine eigenvalue to be one.}
Here an even more useful conclusion holds. Let $V$ be an additive
group, $L\subset V$ a subgroup, and $f:V\to K$ be $L$-periodic,
with only finitely many nonzero $L$-cosets and nonzero sum over them.
Suppose $T$ is an additive automorphism and
\[
 f(Tv+b)=\chi f(v)\quad(v\in V).
 \tag{CI.5}
\]
Then $\chi=1$, even if $T$ does not preserve the initially chosen
period subgroup $L$.

First $f\ne0$ and the surjectivity of $T+b$ give $\chi\ne0$.
Let $M=\{u:f(v+u)=f(v)\ \text{for all }v\}$ be the full period
group. Equation (CI.5) and its inverse imply $T(M)=M$. Fix
$v_0$ with $f(v_0)\ne0$. Distinct classes in $M/L$ give distinct
nonzero support classes $v_0+u+L$, so $[M:L]$ is finite.
Each nonzero $M$-coset contributes exactly $[M:L]$ equal terms to
the sum over $L$-cosets. Hence the sum over $M$-cosets is nonzero.
But $T+b$ permutes these cosets and their finite nonzero support.
Summing (CI.5) gives $S=\chi S$ with $S\ne0$, proving the claim.
This is the actual finite-sum rigidity used in the projected-kernel
argument; root-of-unity containment alone would not prove it.

\paragraph{CI3: finitely many admitted local level cosets.}
Let $\Lambda=\bigoplus_l Oq^{e_l}e_l$, and call $h$ admitted at
depth $j$ if
\[
 h^{-1}\Lambda\subseteq\Lambda,
 \qquad q^j h\Lambda\subseteq\Lambda.
 \tag{CI.6}
\]
Color $h$ by the matrix of $q^jh$ in this lattice basis modulo
$q^{j+1}$. There are finitely many colors. More literally, one may
lift every local integral entry to $\mathcal O_E$ modulo
$q^{j+1}O$ and use its class in the finite global quotient
$\mathcal O_E/(q^{j+1})$, exactly as the source does. Such lifts
exist because a denominator outside $P$ is invertible modulo a
sufficient power of $P$; alternatively localization of this finite
quotient is a quotient of it.

Equal colors give $(h'-h)\Lambda\subseteq q\Lambda$.
Applying $h^{-1}$ and using (CI.6) gives
$(h^{-1}h'-1)\Lambda\subseteq q\Lambda$; interchanging $h,h'$
gives the same assertion for its inverse. Thus $h^{-1}h'$ is an
integral basis change congruent to one modulo $q$.

Every local group $J_j$ in this construction contains that
principal congruence subgroup. At a twist prime it preserves the
two lattices defining $J_Q$; at a type prime it preserves the
inverse images of the residue flags and fixes their prescribed
residue basis vectors. These are direct consequences of being
the identity on $\Lambda/q\Lambda$. Hence equal colors imply
$h^{-1}h'\in J_j$, proving finiteness of admitted right cosets and
therefore of admitted $R_j\backslash\mathrm{GL}(Y_j)/J_j$
classes. The source's type-prime admission condition requires
stabilizing $\Lambda_j$, and therefore implies (CI.6) automatically.
This proves the actual
\texttt{exists\_finite\_index\_of\_adm\_cosets} supplier.

\paragraph{CI4: a rational normal form with integral projected group.}
At a twist prime let $M=hO^{n+c}$ with
$O^{n+c}\subseteq M\subseteq q^{-j}O^{n+c}$. Put
$W_a=X+Ee_0+\cdots+Ee_a$. The image of $M\cap W_a$ under its
$e_a$-coordinate is a nonzero fractional ideal of $O$, hence
$q^{\nu_a}O$. Choose $v_a\in M\cap W_a$ with that coordinate
$q^{\nu_a}$, and put $u_a=q^{-\nu_a}v_a$.
Successively eliminating the last coordinate proves
\[
 M=(M\cap X)\oplus\bigoplus_{a=0}^{c-1}Oq^{\nu_a}u_a,
 \qquad -j\le\nu_a\le0.
 \tag{CI.7}
\]
The inequalities follow by comparing the same coordinate image
with $O$ and $q^{-j}O$. Choose a unitriangular map which fixes
$X,e_0$ and carries $u_a$ to $e_a$ for $a>0$; it exists by
successive elimination. Write $u_0=e_0+z_0$. Since
$\nu_0\le0$, $u_0\in M$, and since $e_0\in M$, one has
$z_0\in M\cap X$. Now choose $g\in\mathrm{GL}(X)$ carrying
the free full $O$-lattice $M\cap X$ onto $O^n$, and extend it
by the identity on the tail. The resulting $\rho\in R_j$ gives
\[
 \rho M=O^n\oplus Oq^{\nu_0}(e_0+z)
          \oplus\bigoplus_{a=1}^{c-1}Oq^{\nu_a}e_a,
 \qquad z\in O^n.
 \tag{CI.8}
\]
If $c=0$, only the first summand occurs. Formula (CI.8) again lies
between $O^{n+c}$ and $q^{-j}O^{n+c}$, so the normalized matrix
is still admitted. It is in the same double coset, with right
factor equal to one.

The intersection of (CI.8) with $X$ is exactly $O^n$.
If $r\in R_j\cap h'J_jh'^{-1}$, then $r$ preserves $h'O^{n+c}$
and $X$, and its inverse does too. Its restriction to $X$ therefore
preserves $O^n$ in both directions. This proves the actual
$\operatorname{pr}_X C_{h'}\subseteq\mathrm{GL}_n(O)$ clause,
as well as \texttt{TwNF}, \texttt{TwBdd}, and \texttt{TwNormRep}.
No global free basis for an $\mathcal O_E$-lattice has been assumed.

\paragraph{CI5: the unfolding fibers are affine spaces, with the stated phase.}
For a step with source tail $f_0,\ldots,f_c$ and target tail
$e_0,\ldots,e_{c-1}$, the matrix $x(t)=\texttt{xx1}(x,t)$ is
described by its columns:
\[
 \begin{array}{c|c}
 \text{source column}&\text{image}\\ \hline
 X'&x(X')\subseteq X\\
 f_0&e_0\\
 f_a\ (1\le a<c)&e_a+\text{an arbitrary vector in }
             X+Ee_0+\cdots+Ee_{a-1}\\
 f_c&\text{an arbitrary vector in }Y.
 \end{array}
 \tag{CI.9}
\]
The fixed row $f_0\mapsto e_0$ applies only when $c>0$.
When $c=0$, the sole source tail column is $f_0=f_c$ and is
the free last column; the target has no tail. The entries
called arbitrary are the coordinates of $t$.
For $r\in R_j$, with restriction $g$ to $X$, column substitution
gives an invertible affine transformation $T_rt+b_r$ such that
\[
 r x(t)=\texttt{xx1}(gx,T_rt+b_r),\qquad
 \operatorname{chArg}(r x(t))=
       \operatorname{chArg}(x(t))+a(r).
 \tag{CI.10}
\]
The inverse comes from $r^{-1}$. For the phase identity, the
$e_a$-coordinate of the image of $f_{a+1}$ acquires precisely
$(r e_{a+1})_{e_a}$ for $a<c-1$; the last column acquires no
constant. All lower terms have zero $e_a$-coordinate.
Right composition by a source $r'$ has the corresponding formula
with retained variable $xg'$ and phase shift $a(r')$. Indeed
each new source column is itself plus earlier columns, so only
the immediately preceding tail column contributes to the phase.
These calculations prove both left and right affine transports,
including the invertibility and surjectivity on the retained
matrix variable.

\paragraph{CI6: twist-prime test functions, covariance, and uniform bounds.}
Reorder coordinates as the first $k$ coordinates of $X$, then
the $c$ tail coordinates, then the remaining coordinates of $X$.
In these coordinates $J_Q$ consists of the matrices in
$\mathrm{GL}(O)$ whose reduction preserves the first $r=k+j+c$
coordinates. Thus its lower-left residue block is zero.
The actual $\Phi_Q$ is supported on integral rectangular matrices
with that lower-left block zero modulo $q$, and on this support is
\[
 \Phi_Q(y)=\xi\bigl(\det\overline{y}_{[0,r),[0,r)}\bigr),
 \tag{CI.11}
\]
where a multiplicative character is extended by zero at zero.
An element $\kappa\in J_Q$ and its inverse preserve this support.
Block multiplication gives
$\overline{(\kappa^{-1}y)}_{11}
=\overline{\kappa^{-1}}_{11}\bar y_{11}$ and hence
\[
 \Phi_Q(\kappa^{-1}y)=
   \xi(\det\overline{\kappa^{-1}}_{11})\Phi_Q(y).
 \tag{CI.12}
\]
The scalar is nonzero, since the leading block and its inverse
are inverse modulo $q$. Right multiplication has the same
determinant calculation. Factoring the full determinant into
the leading and complementary diagonal blocks identifies these
scalars with the exact \texttt{lamRow} and \texttt{muAt} factors
used by the source. In particular \texttt{lamRow} is a unit-valued
character of $J_Q$, and is one on its principal congruence subgroup.

Write the translated datum as
$G_{a,h}(y)=s(a,h)\Phi_Q(h^{-1}ya)$. Its scalar $s(a,h)$ is the
displayed product of local determinant characters and square-root
powers in \texttt{Loc.om}. It is nonzero when the prescribed
uniformizer character value and square root are nonzero:
every nonzero element is uniquely $q^v u$ with $u\in O^\times$,
so its local character is a nonzero power times $\xi(\bar u)$.
If $a$ is admitted at step $i$ and $h$ at step $i+1$, then
\[
 G_{a,h}(y)\ne0\Longrightarrow q^{i+1}y\text{ is integral},
 \qquad G_{a,h}(y+q^{i+1}z)=G_{a,h}(y)
 \tag{CI.13}
\]
for integral $z$. For support, use $y=h(h^{-1}ya)a^{-1}$.
For periodicity, $h^{-1}q^{i+1}za$ is divisible by $q$, because
$h^{-1}$ is integral and $q^ia$ is integral. This preserves every
residue entry in (CI.11), including the zero-support cases.
The phase $\operatorname{chArg}$ is a sum of entries, so its value
on this support lies in $q^{-(i+1)}O$. These are uniform bounds,
not assumptions about the chosen matrix $y$.

\paragraph{CI7: actual twisted double-coset nonvanishing.}
Suppose $h'=r h\kappa$ with $r\in R_j$, $\kappa\in J_Q$, and
$h,h'$ admitted. By (CI.12) and the determinant factors,
\[
 G_{a,h'}(y)=C\,G_{a,h}(r^{-1}y),\qquad C\ne0.
 \tag{CI.14}
\]
Assume the depth-$D$ character $\psi$ is multiplicative under
addition on $q^{-D}O$ and is one on $O$. It is nonzero throughout
$q^{-D}O$, since $\psi(u)\psi(-u)=1$. Take $D\ge j$.
If the translated test function has empty support, both
unfoldings vanish. Otherwise choose a support point.
Where either side of (CI.14) is nonzero, CI6 bounds both phase
arguments in $q^{-D}O$. Their difference, the constant in
(CI.10), is therefore in that group too. Consequently the two
unfolding integrands are related by an invertible affine change
of $t$ and a nonzero scalar, including their zero-support terms.
CI1 proves that their stable sums vanish simultaneously. The
invertible retained-variable change $x\mapsto g^{-1}x$ then gives
\[
 \texttt{FiniteSumNonzeroLargeDepth}(a,h)
 \quad\Longleftrightarrow\quad
 \texttt{FiniteSumNonzeroLargeDepth}(a,h').
 \tag{CI.15}
\]
The same reasoning gives source-side transport whenever the
source covariance law and corresponding support bounds hold.
It uses only finite sums over commensurable lattices. This proves
the twist-prime content of \texttt{TwistedEisensteinNonvanishing}
and \texttt{unfoldingNonvanishingDoubleCoset\_holds}, including
the actual normalized representative from CI4. The source's
analytic-sounding name does not add an analytic hypothesis here.

\paragraph{CI8: the actual projected-kernel character is trivial.}
Put $C_h=R_j\cap hJ_jh^{-1}$ and
\[
 \chi_h(r)=\psi(a(r))^{-1}\lambda_j(h^{-1}rh).
 \tag{CI.16}
\]
At a twist prime, $q^j r$ is integral for $r\in C_h$: use
$r=h\kappa h^{-1}$ and admission of $h$. Thus $a(r)\in q^{-j}O$.
For $D\ge j$, (CI.2), the character law of CI6, and additivity
of $\psi$ prove that (CI.16) is a nonzero multiplicative character.

Suppose a nonzero unfolding enters $h$ from some admitted $a$.
If $r\in C_h$ acts identically on $X_j$, then it is triangular
with every diagonal entry one, so $\det r=1$. The target
covariance formula therefore gives
$G_{a,h}(r y)=\lambda_j(h^{-1}rh)G_{a,h}(y)$ without an extra
determinant factor. Fix a retained $x$ at which the stable
unfolding sum is nonzero. By CI5 its integrand satisfies
\[
 f_x(T_rt+b_r)=\chi_h(r)f_x(t).
 \tag{CI.17}
\]
The support and period bounds of CI6 make it a finite sum on a
quotient of the free-coordinate lattice. CI2 proves
$\chi_h(r)=1$. Thus \texttt{ChiTriv} is proved for the actual
kernel, not merely for an assumed quotient character.

Define $\bar\chi_h(g)=\chi_h(r)$ using any lift of
$g\in\operatorname{pr}_X C_h$. Two lifts differ by the kernel,
so this is well defined and multiplicative. CI4 places the group
inside $\mathrm{GL}_n(O)$. This is the precise \texttt{TwChrOk}
and \texttt{characterNonvanishing\_holds} conclusion. At the last
index $c_j=0$ the projection has trivial kernel. At index zero
$h=1$, the twist character is one: $a(r)\in O$ and the
complementary determinant block in \texttt{lamQB} is empty,
since $n_0=k$.

\paragraph{CI9: recursive finite index sets with exact quantifiers.}
Fix one local prime. Start with $F_0=\{1\}$. Given finite $F_i$,
consider admitted matrices $h$ at step $i+1$ for which
\emph{some} $a\in F_i$ has nonzero unfolding into $h$. CI3 gives
finitely many possible double cosets. At a twist prime choose
the normalized representative of CI4 in each retained class;
CI7 preserves its nonzero predecessor. At a type prime retain
any representative in the class. Choose exactly one per class.
This constructs $F_{i+1}$, and proves, separately, its predecessor
property, uniqueness modulo double cosets, and coverage of every
admitted matrix having a nonzero predecessor. The construction
is for $i<m-2$. Set $F_{m-1}=\{1\}$ and $F_j=\varnothing$ for
$j>m-1$, exactly as in the source.

At the finitely many local primes take
$\beta_j=\prod_P F_{P,j}$. For a prescribed component in a
nonempty $\beta_{i+1}$, predecessor choices at all primes form a
member of $\beta_i$; hence the source's predecessor quantifier
is satisfied. Similarly choose the covering representative
componentwise. This proves the actual \texttt{exists\_step} and
\texttt{exists\_sets} construction. Intermediate index sets are
not asserted nonempty merely because the initial set is nonempty.
Their exact nonemptiness requirements are handled by the
nonvanishing arguments, not by changing an empty product silently.

\paragraph{CI10: the actual symmetric residue pattern and its characters.}
Put $s=m/2-1$, $r=k+m-1$, and use the reordered coordinates of
CI6. For a square matrix of size $N=r+jt$, the allowed off-diagonal
positions $(p,p')$ are precisely
\[
\begin{array}{c|l}
 \text{column }p'&\text{allowed row }p\\ \hline
 p'<k&p<p'\ \text{or}\ k\le p<k+s\\
 k\le p'<r&k\le p<p'\ \text{or}\ (p'>k+s\ \text{and}\ p<k)\\
 p'\ge r&p<r\ \text{or}\
       \lfloor(p-r)/t\rfloor\le\lfloor(p'-r)/t\rfloor.
\end{array}
\tag{CI.18}
\]
Only indices within the matrix are used. When $t=0$ the last row
of the table has no columns. Let $\mathcal A_N$ be the matrices
whose off-diagonal support satisfies (CI.18). The relation
``equal or allowed'' is transitive: for a column in the first
two intervals, its allowed predecessors stay in those intervals
and a second predecessor satisfies the same inequality; for a
column in the last interval, block numbers can only decrease.
The one missing corner, from column $k+s$ to a row $p<k$, cannot
arise through an intermediate index, since that column has only
earlier tail predecessors. Thus $\mathcal A_N$ is an algebra.

On the first $r$ indices the strict allowed relation has no
cycles. For example it strictly increases the ordering function
which first lists $k,\ldots,k+s-1$, then $0,\ldots,k-1$, then
$k+s,\ldots,r-1$. Hence
$(AB)_{pp}=A_{pp}B_{pp}$ for $p<r$ in this algebra. An invertible
$A\in\mathcal A_N$ has inverse in the algebra: multiplication
by $A$ is injective on the finite-dimensional algebra, hence
surjective, so its inverse is the element mapping to one.
It follows that
\[
 J_N^{\rm res}=\{A\in\mathcal A_N^\times:A_{pp}=1\ (p<r)\}
 \tag{CI.19}
\]
is a group. Conjugating the rational matrices by the diagonal
basis of $\Lambda_j$ and reducing modulo $q$, the defining
\texttt{JGsym} flag and residue-vector clauses are exactly
(CI.19). In fact each flag is the inverse image of its span of
earlier reordered basis vectors. Its exceptional boundary
\texttt{bd5} imposes precisely the missing corner just described.
This proves both directions of the entry description.

Define
\[
 \ell(A)=\sum_{a=0}^{k-2}A_{a,a+1}
          +\sum_{a=k}^{r-2}A_{a,a+1},\qquad
 \lambda_j(A)=\psi_0(\ell(A))
       \prod_{b=0}^{j-1}\chi_b(\det A_{[r+bt,r+(b+1)t)}).
 \tag{CI.20}
\]
For each displayed adjacent pair, the matrix product has only
the two possible intermediate indices equal to its endpoints:
the table rules out all others. The omitted pair $(k-1,k)$ is
zero throughout the pattern. Thus $\ell(AB)=\ell(A)+\ell(B)$.
Each diagonal $t$-block multiplies as a block, since all
intermediate indices outside it contribute zero. Its determinant
is nonzero for a group member, by the same statement for its
inverse. Therefore (CI.20) is a unit-valued character and is one
on a matrix reducing to one. This reconstructs the actual
symmetric \texttt{lamRow} character law.

The rectangular test function has $a=r+(j-1)t$ columns and
$a+t$ rows. On its pattern support it is
\[
 \Phi_S(Y)=\psi_0\!\left(\sum_{b=0}^{r-2}Y_{b,b+1}\right)
 \frac{\prod_{b=0}^{j-2}\chi_b(\det Y_{[r+bt,r+(b+1)t)})}
      {\chi_{j-1}(\det Y_{[0,a)})},
 \tag{CI.21}
\]
with the same zero extension when a determinant vanishes.
The pattern is stable under left and right multiplication by
the corresponding groups (CI.19). To verify its covariance,
multiply blocks: each old diagonal block determinant gains the
corresponding determinant of the acting matrix, the leading
$a$-determinant gains its leading determinant, and the adjacent
entry sum gains $\ell$ of that matrix. For target multiplication
the full determinant is the leading $a$-determinant times the
new last $t$-block determinant. These identities give, with the
source's exact local determinant character $\mu$,
\[
 \Phi(\kappa y)\mu(\det\kappa)
    =\lambda_j(\kappa)\Phi(y),\qquad
 \Phi(y\kappa')\mu(\det\kappa')
    =\lambda_{j-1}(\kappa')\Phi(y).
 \tag{CI.22}
\]
The identities also hold off support, because the inverse
matrices preserve support. Reduction of the scaled rational
entries gives the literal \texttt{PhiGS} source and target
laws. This supplies the type-prime covariance needed below,
rather than assuming it as part of a later column-law record.

\paragraph{CI11: a nonzero residue unfolding has the identity double coset.}
In CI11--CI13 assume $j+1<m$, hence $c=c_j\ge1$; this is
the source's one-coset range. The zero-tail projected kernel
is handled separately in CI8.
Let $h\in\mathrm{GL}_{a+t}(\mathbb F_q)$. If the finite
unfolding of $\Phi_S(h^{-1}\texttt{xx1}(x,t))$ is nonzero,
some summand has $Y=h^{-1}\texttt{xx1}(x,t)$ on nonzero support.
Its leading $a\times a$ block is invertible by (CI.21).
Complete the $a$ columns of $Y$ by the last $t$ standard columns.
The resulting square matrix $\kappa$ is invertible, has determinant
that leading determinant, and belongs to $J_{a+t}^{\rm res}$
by (CI.18).

From (CI.9) choose $\nu\in N_j$ whose tail columns are the first
$c$ tail columns of $\texttt{xx1}(x,t)$; its triangular diagonal
is one. Then $g=\nu^{-1}h\kappa$ fixes every tail basis vector.
Moreover it sends into $X$ each reordered source column outside
the tail interval and its one free last column. Let $P_X,P_A$
be the projections onto $X$ and the tail, and put
$B=P_AgP_X$. Since $B$ kills the tail and takes values in it,
$B^2=0$. The matrix $\kappa'=1+B$ is invertible with inverse
$1-B$. The only possible nonzero entries of $B$ have a tail row
and a column at or after the omitted free column; all these
positions are allowed by (CI.18). Thus $\kappa'\in J_{a+t}^{\rm res}$.
For $\gamma=g(1-B)$ the tail component of $\gamma v$ equals that
of $v$, while $\gamma$ fixes the tail and maps $X$ into $X$.
Consequently $\gamma\in R_j$. We obtain
\[
 h=\nu\gamma\kappa'\kappa^{-1}\in R_jJ_{a+t}^{\rm res}.
 \tag{CI.23}
\]
This proves the support part of \texttt{OneCosetSharp} by matrix
completion and elimination; cancellation in a Fourier sum has
not been assumed to identify its support.

\paragraph{CI12: an explicit nonzero residue unfolding at the identity.}
Choose $x_0:E^{n'}\to E^n$ on coordinate vectors by
$x_0e_b=e_b$ for $b<k$ and $x_0e_b=e_{b+1}$ for $b\ge k$.
Use this same matrix over the residue field. For any free
coordinates $t$, the reordered columns of
$Y=\texttt{xx1}(x_0,t)$ outside $k<p'\le k+c$ are standard
columns. If $Y$ satisfies the support conditions of (CI.21), its
leading block is upper triangular with diagonal one, and its
old $t$-blocks are identity matrices. Indeed in the remaining
columns (CI.18) permits only rows strictly earlier than the
column; all other columns were already standard. Hence all
multiplicative-character factors in (CI.21) are one.
The only possibly nonzero terms of its adjacent-entry sum are
the $c$ terms of $\operatorname{chArg}(Y)$. Therefore
\[
 \Phi_S(\texttt{xx1}(x_0,t))
       \psi_0(-\operatorname{chArg}(\texttt{xx1}(x_0,t)))
   =\begin{cases}1&\text{if the support pattern holds},\\
                  0&\text{otherwise}.
     \end{cases}
 \tag{CI.24}
\]
There is at least one supporting choice: all free entries zero
except the entry sending the last source tail vector $f_c$ to
the target vector $e_k$ in $X$, which is one. In reordered
coordinates this is the rectangular identity matrix. The finite
unfolding at $x_0$ is thus the positive integer counting these
supporting choices. It is nonzero in characteristic zero.
This proves the nonvanishing half of \texttt{OneCosetSharp};
even a nontriviality hypothesis on $\psi_0$ is unnecessary for
this particular finite-matrix calculation.

\paragraph{CI13: reduction, lifting, and the actual type-prime one-coset theorem.}
Let $h$ stabilize $\Lambda_j$. In its scaled basis, both $h$ and
$h^{-1}$ are integral and reduce to inverse residue matrices.
Its determinant is a unit, so the scalar in the translated
$\Phi_{GS}$ is nonzero. The scaled entries of an unfolding
matrix are again of the form (CI.9). More precisely, for a free
entry in row $l$ and source column $b$, put
\[
 d_f=(b-n'+1)-\begin{cases}0&l<n,\\l-n+1&l\ge n.
                       \end{cases}
 \quad 1\le d_f\le c+1.
 \tag{CI.25}
\]
Its scaled entry is $q^{d_f}t_f$. The retained $x$ entries
remain unscaled. If one is nonintegral the unfolding is zero,
since every nonzero test-function value requires integral
scaled entries.

For integral $x$, reduce it and every $q^{d_f}t_f$ modulo $q$.
In a box (CI.3) with $A\ge c+1$, $B\ge1$, every residue value
has exactly $q^{B+d_f-1}$ admissible representatives in the
$f$th free coordinate. This follows by writing its numerator
as $q^{A-d_f}(u+qv)$, where $u$ is the chosen residue digit.
The phase satisfies
$\overline{q\operatorname{chArg}(\texttt{xx1}(x,t))}
=\operatorname{chArg}(\texttt{xx1}(\bar x,\bar t))$.
The literal type-prime tie
$\psi(a/q)=\psi_0(\bar a)$ consequently gives the exact identity
\[
 \sigma(E^{A,B}_{1,h}(x))
   =\sigma(s(1,h))\,
        q^{\sum_f(d_f-1)} E^{\rm res}_{1,\bar h}(\bar x).
 \tag{CI.26}
\]
Here the residue function is (CI.21), and $\sigma:K\hookrightarrow
\mathbb C$ is the specified embedding. All factors in front are
nonzero. Every residue $x$ lifts entrywise to an integral rational
$x$. Thus eventual nonzero unfolding is equivalent to nonzero
residue unfolding, in both directions and with the stated
uniform depths.

Apply CI11 to $\bar h=\bar\rho\bar\kappa$. Lift the entries of
$\bar\rho$ to $O$, preserving its zero entries and its entries
equal to one. Its leading $X$ block has unit determinant; its
triangular tail has diagonal one. Hence the lift and its inverse
are integral. Conjugating back by the scaled lattice basis gives
$\rho\in R_j\cap\operatorname{Stab}(\Lambda_j)$ with that
reduction. The matrix $\rho^{-1}h$ has residue $\bar\kappa$ and
therefore belongs to $J_j$ by CI10. Consequently
\[
 E_{1,h}\ne0\ \Longrightarrow\ h\in R_jJ_j,
 \qquad E_{1,1}\ne0.
 \tag{CI.27}
\]
The second assertion follows from CI12 and (CI.26).

For completeness these statements transport in both source and
target double cosets, with the stabilizer hypotheses present
in the source. All source and target matrices then preserve the
scaled lattices, so their moved Hom lattice is the same; it has
bounded coordinates and a full period lattice. In that Hom
lattice $q\operatorname{chArg}$ is integral, because successive
tail exponents differ by one. Formula (CI.22), CI5, and the finite
affine change of CI1 prove the two-sided transport exactly as in
CI7, now with depth one. In particular, if the source matrix
$a$ is in the identity double coset and both matrices stabilize
their prescribed lattices, then
\[
 E_{a,h}\ne0\quad\Longleftrightarrow\quad h\in R_jJ_j.
 \tag{CI.28}
\]
This is the actual auxiliary-prime input to \texttt{twSurjG},
\texttt{twNonVanG}, and the reference-chart index matching.
It has been proved independently of \texttt{TC.Laws}.

\paragraph{CI14: the projected symmetric levels and their exact characters.}
At the identity representative, restricting $r\in R_j\cap J_j$
to $X_j$ gives an element of $\mathrm{GL}_{n_j}(O)$, since
$\Lambda_j\cap X_j=O^{n_j}$. Its reduction $A$ has diagonal one
for $p<k+j$; its allowed off-diagonal positions are
\[
\begin{array}{c|l}
 p'<k&p<p'\ \text{or}\ (k\le p\ \text{and}\ p+s+2\le k+j)\\
 k\le p'<k+j&k\le p<p'\ \text{or}\
                     (p<k\ \text{and}\ k+j<p'+s+1)\\
 p'\ge k+j&p<k+j\ \text{or}\
       \lfloor(p-(k+j))/t\rfloor\le
                 \lfloor(p'-(k+j))/t\rfloor.
\end{array}
\tag{CI.29}
\]
Indeed deleting the $c_j$ tail positions from (CI.18) replaces
each remaining index $p\ge k$ by $p+c_j$, and
$c_j+j=2s+1$. Substitution in the three inequalities gives
exactly (CI.29). Conversely lift a residue matrix satisfying
(CI.29) entrywise, with determinant a local unit, and extend it
by the identity on the tail. The same substitutions show that
this extension belongs to $R_j\cap J_j$. This proves the exact
projected group, including surjectivity onto the residue group.

For such an extension of $g$ with residue $A$, split the
superdiagonal sum (CI.20) into the initial $X$ part, the tail
part, the tail-to-$X$ boundary, and the remaining $X$ part.
The boundary entry is zero because $rX\subseteq X$; the tail
part is $q a(r)$ in scaled coordinates. The old block determinants
are precisely the $X$ block determinants. Using
$\psi(a(r))=\psi_0(\overline{q a(r)})$, (CI.16) therefore becomes
\[
 \bar\chi_j(A)=
 \psi_0\!\left(\sum_{b=0}^{k-2}A_{b,b+1}
             +\sum_{b=0}^{j-2}A_{k+b,k+b+1}\right)
 \prod_{b=0}^{j-1}
       \chi_b(\det A_{[k+j+bt,k+j+(b+1)t)}).
 \tag{CI.30}
\]
The two copies of the tail phase cancel, with the inverse
on $\psi$ exactly as in (CI.16). Formulas (CI.29)--(CI.30) are
the source's \texttt{colLevSym} group and character; at $j=0$
they give \texttt{colLev Ch 0}. Thus the actual
\texttt{ocbGrpGLoc}, \texttt{ocbChrGLoc}, and
\texttt{closedLevels\_by\_entries} assertions are proved.

The same calculation also makes double-coset character
transport explicit. If $h'=\rho h\kappa$, then
$C_{h'}=\rho C_h\rho^{-1}$. Its projected group is conjugated
by $g=\rho|_X$. By (CI.2), $a(\rho r\rho^{-1})=a(r)$, while
$\lambda(\kappa^{-1}(h^{-1}rh)\kappa)=\lambda(h^{-1}rh)$.
Hence the projected characters agree under that conjugation.
At type primes the matrices involved preserve $\Lambda_j$, so
$g\in\mathrm{GL}_{n_j}(O)$ and this is also the actual residue
level conjugation used in reference-chart matching.

\paragraph{CI15: constructing the full initial local-index record.}
Assume the literal inputs of \texttt{exists\_chIdx\_of\_cst}:
every row prime is degree one and unramified; twist and type
prime sets are disjoint; $\texttt{sharp}=\texttt{true}$;
\texttt{LocCst.Ok}, $m\le M\le A,B$, the specified unit $\xi_0$
and Weil tie, the exact local determinant-character values and
their type-prime residue tie, and
\[
 D_P=1\quad(P\text{ type}),\qquad
 D_P\ge2M,\quad D_P\ge2(m-1)+1\quad(P\text{ twist}).
 \tag{CI.31}
\]
These are explicit arithmetic and constant hypotheses, rather
than an already supplied index record. CI9 constructs its finite
sets and representatives. At type primes all chosen representatives
are in the identity double coset: induct from $F_0=\{1\}$ using
CI13. That local set is nonempty at every constructed step,
because transport of $E_{1,1}\ne0$ supplies an eligible next
representative. CI14 therefore constructs its residue level and
character by conjugating the identity level. At twist primes use
CI8's projected group and descended character. At zero choose
the literal top twist level and \texttt{colLev Ch 0}; at the last
step the zero tail makes kernel descent immediate.

Here are the remaining bounds required by the record. They are
included to avoid concealing them in a record-assembly assertion.
For type primes write $\mathcal H_i=\operatorname{Hom}_O
(\Lambda_i,\Lambda_{i+1})$. The datum is supported in
$\mathcal H_i$ and invariant under $q\mathcal H_i$, by its
integral-entry and residue definition. Admitted source and target
matrices preserve this same Hom lattice. Its coordinate exponents
are differences of numbers between $0$ and $m-1$; hence its
support is contained in $q^{-m}$ times the ordinary integral Hom
lattice, and its periods contain $q^m$ times that lattice.
Also $q\operatorname{chArg}$ is integral on its support.
At twist primes CI6 gives the sharper bound $i+1\le m$.
Multiplication by the phase preserves these uniform period
bounds, because the change of phase argument is integral.
CI1 now proves the literal \texttt{deep} equality for all
$A'\ge A$, $B'\ge B$.

The retained $X$ entries of \texttt{xx1} are actual matrix
entries. A nonzero summand therefore bounds them by $q^{-A}$;
adding a matrix in $q^B\operatorname{Hom}(O^{n_i},O^{n_{i+1}})$
changes every summand by a period and leaves its phase unchanged.
Multiplication by the fixed \texttt{volN} factor in
\texttt{XiAt} preserves both statements. At the terminal
\texttt{joinCol} step the source tail has one coordinate and
the target tail is zero. In the scaled bases its extra column
is $q y$. Thus at a type prime support requires $q y$ integral,
and $y\mapsto y+z$ with $z\in O^n$ changes that column by a
$q$-period. The $x_0$ columns are integral and have $q$-periods.
At a twist prime the same argument with (CI.13) gives support
$q^{-M}$ and periods $q^M$ for $y$, and the weaker $q^{-A}$
support and $q^B$ periods for $x_0$. These prove exactly
\texttt{psiSupp}, \texttt{psiInv}, \texttt{xiSupp},
\texttt{xiInv}, \texttt{psiSuppX}, and \texttt{psiInvX}.

The projected integrality field follows from CI4 at twist
primes and $\Lambda_j\cap X_j=O^{n_j}$ at type primes.
The four group/character fields follow from the constructions
above: at type primes a residue matrix and any integral lift
give the same value by (CI.30), transported by CI14; at twist
primes any two projected lifts give the same value by CI8.
The zero and last uniqueness and matrix fields, predecessor,
injectivity, and coverage fields are precisely CI9. The
arithmetic character, depth, and Weil fields are the stated
inputs (CI.31). This checks every field of
\texttt{ChIdx.Ok} and proves the actual
\texttt{exists\_chIdx\_of\_cst}, including its twist normal-form
conclusion. No \texttt{Forms0}, \texttt{TC.Laws}, theta-algebraicity
or \texttt{Later} record has entered this proof.

\paragraph{CI16: the actual local constants, including the auxiliary Gauss sum.}
The constant hypotheses used in CI15 can themselves be
constructed from the root and local-character data already
isolated in TL14--TL15. Suppose $\xi_0\in O^\times$, and the
coefficient field contains a primitive $q^D$th root. Adjust its
power so that its specified complex embedding is
$\zeta=\exp(2\pi\mathrm i/q^D)$; two primitive roots differ by
a power coprime to $q^D$. For $x\in q^{-D}O$, let $b$ satisfy
$\xi_0x-b/q^D\in O$, and define
\[
 \psi(x)=\zeta_K^b;
 \qquad\psi(x)=1\text{ outside }q^{-D}O.
 \tag{CI.32}
\]
CI1 gives $b$, uniquely modulo $q^D$. Therefore the definition
is independent of its representative, is additive into
multiplication on $q^{-D}O$, is one on $O$, and is not one at
$q^{-1}$, since $\xi_0$ is a local unit. The stated Weil formula
holds for every allowed integer representative $b$, not just
the chosen one. At a type prime take $D=1$; the supplied
\texttt{Ch.WeilTied} formula identifies (CI.32) with
$\psi_0(\overline{qx})$. Thus its \texttt{psi0} clause is also exact.

For the canonical type-prime square root the source uses the
quadratic Gauss sum, not an arbitrary root. Let $\epsilon$ be
the quadratic character of $\mathbb F_q^\times$, extended by
zero, and $G=\sum_x\epsilon(x)\psi_0(x)$. Here $q\equiv1\pmod4$
and $\psi_0\ne1$. Translation of the complete additive sum
shows $\sum_x\psi_0(ax)=0$ for $a\ne0$, and multiplication by
a nonsquare shows $\sum_x\epsilon(x)=0$. With $y=xt$ in $G^2$,
\[
 G^2=\sum_t\epsilon(t)\sum_{x\ne0}\psi_0(x(1+t))
     =\epsilon(-1)(q-1)-\sum_{t\ne-1}\epsilon(t)=q.
 \tag{CI.33}
\]
The last equality uses $\epsilon(-1)=1$, since $4\mid q-1$.
This proves the actual \texttt{ciSq} square identity over $K$.
At twist primes a supplied square root witnesses the predicate
defining \texttt{twSqrt}. Its chosen value therefore squares to
$q$ and is nonzero. No sign or positive-embedding normalization
is part of that definition.

Choose $\mu_q$ to be the specified inverse reference local
character value times the inverse twist value at type primes;
at twist primes use the prescribed product of the two inverse
coefficient-field units. Their nonvanishing and all \texttt{muq}
equalities follow directly from those values and injectivity
of the embedding. The residue-unit tie is the literal
\texttt{Ch.Tied} clause. This constructs
\texttt{exists\_locCst\_of\_values}.

The tied data with integral weights also supply the required
reference values at ramified primes. Here is the additional
argument behind \texttt{locK\_mem\_of\_tied}. For a fixed local
element $x\in E^\times$ at $P$, choose $a\in E^\times$ with
$a/x\equiv1\pmod P$ and $a\equiv1$ at the other finitely many
ramified primes. This is elementary approximation: clear the
denominators by an integer $b$, solve for an algebraic integer
$y$ modulo the requisite powers of the distinct prime ideals
by the Chinese remainder theorem, and set $a=y/b$; choose a
nonzero solution. Tameness gives $\mu_P(a)=\mu_P(x)$ and makes
the other ramified factors one. At every unramified prime the
local value at $a$ is an integral power of its tied uniformizer
value in $K$, and only finitely many such powers are nonzero.
The infinity type has exponents $-2z_\tau$ with $z_\tau\in\mathbb Z$,
so its value is a product of
$(\tau(a)/\overline{\tau(a)})^{-z_\tau}$, also in the embedded
coefficient field. The Hecke product formula now places
$\mu_P(x)$ in that field. This proves the reference-value
supplier from the literal tied integral-weight data, and gives
the full \texttt{exists\_locCst} used for the initial indices.

For initial indices take $M=A=B=m$, $D_Q=2m$ and
$D_G=1$. For canonical indices the actual definition is
$M=m$, $A=B=m+1$, $D_Q=2m$, $D_G=1$, with
\texttt{sharp=true}. The nine fields of \texttt{LocCst.Ok}
now follow from (CI.32), (CI.33), the supplied twist square
roots, and these inequalities. In particular the larger canonical
box depths do not change $M$ to $m+1$.

\paragraph{CI17: the specific canonical indices satisfy their full record.}
Consider the source's actual \texttt{twCst} and \texttt{twIdx},
not an arbitrary existence witness. Assume $X.\texttt{Ok}$,
$Pr.\texttt{Ok}$, $Ch.\texttt{Ok}$, the given character and
Weil ties, type-prime inverse reference values in $K$,
$q_P\equiv1\pmod4$ at type primes, degree-one unramified twist
primes disjoint from them, and the root, square-root, and
local-unit clauses of \texttt{TwCstQ}.

The canonical step selects a class at each twist prime among
admitted matrices having a nonzero predecessor in the preceding
global index set. CI3 makes this class set finite, so its
``finite if finite, empty otherwise'' definition takes the
finite branch. CI4 and CI7 provide a member of every such class
satisfying all predicates in the choice of \texttt{twRep}:
admission, normal form, projected integrality, and the original
nonzero-predecessor condition. Thus the chosen representative
has those properties, regardless of which eligible witness
the choice operation returns.

The canonical matrix is one at every type prime. The definition
of \texttt{twNext} retains its explicit condition that every
type prime have a predecessor into that identity matrix.
CI13 proves the identity nonzero assertion and the one-coset
support statement needed for this condition. For coverage,
start with any tuple of admitted matrices having the required
predecessors: its twist classes choose the product component,
and its type components are in the identity double coset by
CI13. Transport to the identity supplies the defining condition
of \texttt{twNext}. Conversely this condition and the twist
candidate condition give every requested predecessor. Equality
of all local double cosets gives equality of the selected
product of classes, hence equality of indices. This proves the
actual predecessor, injectivity and coverage fields, without
assuming that a candidate set is nonempty.

At positive twist steps CI8 proves the actual character used
by \texttt{twLam}; at the last step its projection kernel is
zero, and at zero it is the trivial top level. At all type
steps CI14 identifies the group and character with the actual
\texttt{colLevSym} definition. The support, periodicity and
depth fields follow from the uniform calculations in CI15,
now with the canonical constants of CI16. The base and terminal
index types are the one-element type stipulated by the source,
and the types above the terminal step are empty. The arithmetic
ties come from CI16 and the present literal hypotheses.
This proves \texttt{twIdx\_ok\_closed}; its formerly separate
\texttt{TwFinQ}, \texttt{TwNormDC}, \texttt{TwChrQ},
\texttt{TwZeroQ}, and \texttt{TwAnQ} suppliers have all been
proved by CI1--CI16.

\paragraph{CI18: canonical replacement really preserves \texttt{Ok}.}
Let $I$ be an actual supplied column datum with
$I.\texttt{Ok}$ and old twist depths at least $2m$, as in the
source's \texttt{Deep} hypothesis. TL20 already extracted its
primitive $q^{2m}$th roots, square roots, and inverse reference
local values. To recall the root step precisely: if the old
depth is $D=2m+b$, insert
$x=\xi_0^{-1}/q^{2m}$ and integer $q^b$ in its Weil identity.
The local-unit clause verifies both allowed denominators, and
the resulting value has complex image
$\exp(2\pi\mathrm i/q^{2m})$. Its order is exactly $q^{2m}$
by injectivity. This extracts the arithmetic premises only;
the nonvanishing and kernel arguments are supplied by CI7,
CI8, and CI10--CI14, not by this extraction.

Apply CI17 for each tuple in $I$. Replacing its constants and
indices by the defined \texttt{twCst} and \texttt{twIdx} gives
their actual \texttt{ChIdx.Ok}. Every other field of $I$ is
unchanged. The coordinate, host, character, tuple-reading,
selection, and ambient-data fields of \texttt{TC.Ok} are
therefore inherited literally. This proves
\texttt{twIdx\_ok\_of\_ok} and \texttt{canon\_ok};
\texttt{IsTw} follows from the definitions. TL13 then derives
the column laws for the canonical datum. No use of those laws
was made in proving its local-index predicates.

\paragraph{CI19: the resulting canonical evaluation construction boundary.}
For the initial constructor, choose the finite character family,
tuples, residue characters and coefficient fields of TL14--TL16.
Supply the literal field/host/selection data and the ambient
period datum of the constructor. CI16 gives local constants
with $M=A=B=m$ and twist depth $2m$; CI15 gives their actual
indices. Installing those records therefore supplies the
local-index fields of the initial \texttt{TC.Ok}. The other
fields are precisely the supplied raw constructor data, and
must be checked by their own arithmetic and ambient suppliers.
CI18 now gives a canonical datum with \texttt{Ok} and
\texttt{IsTw}, and TL13 gives \texttt{Laws}.

Accordingly the local-coset, normal-form, projected-character,
and first-unfolding boundary in TL20--TL21 is discharged. The
remaining hypotheses in that \texttt{EvLg} construction include
the actual host and raw arithmetic data and the normalized-theta
vacuum/indicator identities, smoothness and covariance used by
TL18--TL19 for \texttt{JetsIn}. In this version PA9 separately
supplies the literal ambient period/trace-isometry record,
and PA13 supplies its CM-field conclusion. They can be combined
with the present construction using the same ambient datum;
neither was a premise of CI1--CI18. AH17--AH18 separately construct
the arithmetic host under their exact prime and splitting hypotheses,
and TA10/TA12 supply the one-dimensional normalized theta arithmetic.
These do not by themselves identify the actual vacuum and indicator
family or prove its covariance. The coefficient-field containment
antecedents and their original order of quantifiers remain
unchanged. None of this produces the five independent stages
of \texttt{Later}, or proves RH or GRH from the local index
calculation alone.
